\section{Methods}
\label{sec:proposed}
In this section, we investigate methods to mitigate leakage errors in X and CZ gates, beginning with the X gate. Our
approach builds on the DRAG-1 scheme. Conventionally, DRAG-1 relies on a seed function to generate a control pulse to
minimizes leakage from $\vert{}0\rangle$ to $\vert{}2\rangle$. However, we demonstrate that bypassing the seed function
to synthesize the pulse directly yields a superior pulse shape.
% \subsection{X gate}
\subsection{Note on $Y$ DRAG correction}
Aside from the $X$ control pulse, we utilize the DRAG pulses \ref{eq:ydrag} and \ref{eq:zdrag}.
We demonstrate that the analytical expression for the $Y$ control pulse remains valid even in the non-perturbative
regime---a derivation that, to the best of our knowledge, has not been previously reported in the literature. To
illustrate this, we recalculate the system Hamiltonian without relying on perturbative assumptions. We introduce a
unitary frame transformation $V$, defined as:
% We can show that the formula for the $Y$ control pulse is still valid in non-perturbative regime.
% (Within the knowledge of the author, this computation has not been done/shown elsewhere.)
% Now we recalculate the Hamiltonian without the assumption of perturbation.
% % Applying a frame change with $V=\exp(S)$, $S=-i\lambda_2\dfrac{\Omegax}{2\Delta_2}\hat{\sigma}^y_{1,2}$ to the Hamiltonian \ref{eq:Hrwa} ,
% % the new $\ket{1}\bra{2}$ component reads
% Consider a frame change $V$
\begin{equation}
  V = \exp\left[-i\lambda_2\dfrac{\Omegax}{2\Delta_2}\hat{\sigma}^y_{1,2}\right]
  = \cos(\phi) -i\hat{\sigma}^y_{1,2}\sin(\phi)
  \label{eq:first_frame_change}
\end{equation}
where $\phi = \lambda_2\dfrac{\Omegax}{2\Delta_2}$, and apply it to our qudit Hamiltonian in \eref{eq:Hrwa}.
The $\ket{1}\bra{2}$ component of the resulting Hamiltonian is
$$
  -\sin(2\phi)\dfrac{\delta(t)+\Delta}{2}+\cos(2\phi)\dfrac{\lambda \Omegax}{2}-i\dfrac{\lambda}{2}\left(\dfrac{\Omega_x^\prime(t)}{\Delta}+\Omega_y(t)\right).
$$
% Setting $\Omegay$ as in \eref{eq:ydrag} still reduces this component.
% Hence, we just use the standard DRAG correction \eref{eq:ydrag} for the $Y$ control.
Substituting $\Omega_y$ as defined in \eref{eq:ydrag} demonstrates that this off-diagonal component is still effectively
suppressed. Consequently, we retain the standard DRAG correction  for the $Y$ control pulse.
% \subsection{RWA}



% \subsubsection{Antidiagonalization of $2\times 2$ matrix}
%
% \begin{boxA}
%   {\bf Probelm } \\
%   Given a Hermitian matrix $H\in M_2(\CC)$, antidiagonalize $H$ so that its diagonal elements be minimized according to transformation~(\ref{eq:framechange}).
% \end{boxA}
% Suppose $H$ is a $2\times2$ Hermitian matrix.
% Consider a problem of anti-diagonlizing $H$ according to the transformation~\ref{eq:framechange}.
% Here we mean by anti-diagonalization making the diagonal components as small as possible in the sense of $l_1$ or $l_2$ sum.
%
% $$
%   H    = H_d + H_a, \qquad
%   H_d  = \begin{pmatrix} a(t) & 0 \\ 0 & d(t) \end{pmatrix},\qquad
%   H_a  = \begin{pmatrix} a(t) & 0 \\ 0 & d(t) \end{pmatrix}
% $$
% Consider the following transformation:
% \begin{equation}
%   U(t) = \exp\left(i\int_0^t\,H_d(t^\prime)\mathrm{d}t^\prime\right) = \begin{pmatrix}\exp\left(i\int_0^t\,a(t^\prime)\mathrm{d}t^\prime\right) &0\\0&\exp\left(i\int_0^t\,d(t^\prime)\mathrm{d}t^\prime\right) \end{pmatrix}
%   \label{eq:antidiagonalization}
% \end{equation}
% Since $i\dot{U}U^\dagger=-H_d(t)$, this yields a complete anti-diagonalization
% $$
%   H^\prime(t) =
%   \begin{pmatrix}
%     0                 & b(t)e^{i\phi(t)} \\
%     c(t)e^{-i\phi(t)} & 0                \\
%   \end{pmatrix}
% $$
% where $\phi(t) = \int_0^t\,[a(\tau)-d(\tau)]\mathrm{d}\tau$.
%
% To make this method effective, we need to ensure that $\ket{0}$ and $\ket{1}$ has not changed at the ends of the
% operation by the transformation~\ref{eq:antidiagonalization}. In formula,
% $$
%   \phi(T) \in \RR_{>0}
% $$

\subsection{DRAG-1}
Although in \ref{subsec:drag}, the DRAG-1 pulse is created from the seed pulse waveform $\Omega_1$
through \eref{eq:drag1}, now we reverse this process.
First we solve \eref{eq:drag1} for $\Omega_1(t)$.
We can rewrite  \eref{eq:drag1} as a second-order linear oridnary differential equation (ODE):
\begin{equation}
  \Omega^2_x(t)=\left(1+\frac{D^2}{\Delta_2^2}\right)\Omega_1^2(t).
  \label{eq:drag1_ode}
\end{equation}
Recognizing this as a driven harmonic oscillator equation, the general solution is
\begin{equation}
  \Omega_1^2(t) = \int_0^t\,\left[\dfrac{1}{\Delta_2}\sin(\Delta_2(t-\tau))\right][\Delta_2^2\Omega_x(\tau)^2]\,\dd \tau
  + A\cos(\Delta_2t+\phi).
  \label{eq:drag1_sol}
\end{equation}
The first term represents the particular solution to the ODE in \eref{eq:drag1_ode}.
The verification can be done with
the Leibniz rule \cite{apostol}, which governs the differentiation under the integral sign:
\begin{equation}
  \frac{d}{dt} \int_{a(t)}^{b(t)} F(x,t) dx = F(b(t), t) \cdot b'(t) - F(a(t), t) \cdot a'(t) + \int_{a(t)}^{b(t)} \frac{\partial}{\partial t} F(x,t) dx.
  \label{eq:leibniz}
\end{equation}
Recall we have three boundary conditions on our debt:
\begin{enumerate}
  \item $\Omega_x = 0$ at both ends ($t=0$, $T$): This ensures compatibility with the SWT and enforces boundary continuity for the X-pulse, mitigating spectral leakage.
  \item $\dot{\Omega}_x = 0$ at both ends ($t=0$, $T$): This guarantees a smooth envelope derivative, maintaining essential boundary continuity for the corresponding Y-pulse.
  \item $\Omega_1 = 0$ at both ends ($t=0$, $T$): This further enforces strict SWT compatibility at the start and completion of the gate.
\end{enumerate}

Applying the third boundary condition at the initial time $t = 0$  immediately eliminates the homogeneous component, yielding $A = 0$ in \eref{eq:drag1_sol}
The constraint at $t=T$ translates to
$$
  \left(\Omega_1^2(T)\right) = 0 = \int_0^T\,\left[\dfrac{1}{\Delta_2}\sin(\Delta_2(T-\tau))\right][\Delta_2^2\Omega_x(\tau)^2]\,\dd \tau,
$$
differentiating both sides we get
$$
  \left(\Omega_1(T)\dot{\Omega}_1(T)\right) = 0 = \int_0^T\,\left[\cos(\Delta_2(T-\tau))\right][\Delta_2^2\Omega_x(\tau)^2]\,\dd \tau.
$$
These two constraints combined implies
$$
  \int_{-\infty}^{\infty}[\Delta_2^2\Omega_x(\tau)^2]e^{-i\Delta\tau}\,\dd \tau.
$$
Hence, The boundary condition at $t=T$ for smoothness translates to the null fourier transform condition,
\begin{equation}
  \int_0^T \Omega_x(\tau)^2 e^{-i\Delta_2\tau} d\tau = 0.
  \label{eq:drag1_fourier}
\end{equation}
\subsubsection{Pulse Design}
Although Jesus et al. \cite{99999} proposed a pulse design method for $\Omega_x$ by way of $\Omega_1$, and investigated
an optimal pulse waveform therein, we can
directly design a DRAG-1 pulse through \eref{eq:drag1_fourier}.
Now we consider pushing the gate time limit for maintaining the perturbative approximation of Schrieffer-Wolff transformation.
In it,
The problem of early truncation will be somewhat relieved by keeping $\max |\Omega_x|$ small.
We combine delayed leakage reduction technique described in \cite{Polat2026} and function design techniques to obtain
such pulses.

Let $t_d=\pi/\Delta$. Then $f(t)+f(t-t_d)$'s Fourier transform is computed as
$$
  \mathcal{F}[f(t)+f(t-t_d)](\Delta) = \mathcal{F}[f](\Delta)+e^{i\Delta t_d}\mathcal{F}[f](\Delta) = 0.
$$
This suggests that we find a function $g(t)$ on $[-1,1]$ that satisfies the following property:
\begin{enumerate}
  \item \label{itm:g1} $g$ is anti-symmetric
  \item \label{itm:g2} $g$ is monotonically increasing
  \item \label{itm:g3} $g(1)=1/2$
  \item  \label{itm:g4} $g$'s first and second derivative vanishes at $g(-1)$.
\end{enumerate}
$g$ represents the function on $[0,t_r]$.
On obtaining such a $g$, we construct $\Omega_x^2$ as
$$
  \Omega_x^2(t) =
  \begin{cases}
    A \left(g(2t/t_r-1)-g(-1)\right) \quad                     & (0<t<t_r)       \\
    A                         \quad                            & (t_r<t<t_d)     \\
    A\left( g(2(t_r+t_d-t)/t_r-1)-g(-1)\right)           \quad & (t_d<t<t_d+t_r) \\
  \end{cases}
$$
where $A$ is the amplitude determined by integrating $\Omega_x^2$ according to \eref{eq:xgate_condition}.

There are mutliple choices for such $g$.
\begin{enumerate}
  \item $\dfrac{1}{2} x+\dfrac{1}{2\pi}\sin( \pi x)$
  \item $3\pi x+4\sin(\pi x)+\dfrac{1}{2} \sin(2\pi x)$
\end{enumerate}
\subsubsection{$g$ design}
Let $h$ be a positive symmetric function on $[-1,1]$ that vanishes at both ends.
Put $g^\prime = h^2$. This ensures that the condition \ref{itm:g4} holds as $g^{\prime\prime}=2hh^\prime=0$
and it is easy to check that the other condtions \ref{itm:g1}-\ref{itm:g3} are also satisfied.

\subsection{Optimization by Riccati engineering}
As we've seen in \ref{subsec:npad}, NPAD has an obvious issue that the Hamiltonian is assumed to be
time-independent, and in Ref.~\cite{NPAD2022}, the author only applies this method to perturbative cases, in which the
method reduces to recursive Schrieffer-Wolff transformations.

We outline the newly proposed procedure to handle non-perturbative cases.
% Although well known as a qunatum control method, Ricatti engineering is not widely applied to one qubit gate control.
% In this subsection, we see how Riccati engineering described in \cite{Riccati2010} reduces to in the case of
% dimension $3$.
% Given a $3\times 3$ matrix $H$, we consider the problem of block diagonalizing it into $2\times 2$ and $1\times 1$
% compartments.
% Write
% $$H(t) = \begin{pmatrix} H_11(t) & H_{12}(t) \\ H_{21}(t) & H_{22}(t) \end{pmatrix}$$
% $H'(t) = U(t)H(t)U^\dagger(t) + i\dot{U}(t)U^\dagger(t)$.
% $$U(t) = \begin{pmatrix} (I_2 + X^\dagger X)^{-1/2} & X^\dagger (I_{n-2} + X X^\dagger)^{-1/2} \\ -X (I_2 + X^\dagger X)^{-1/2} & (I_{n-2} + X X^\dagger)^{-1/2} \end{pmatrix}$$
% where $X(t)=(x_1(t)\:\:\: x_2(t))$
% \begin{equation}
%   i\dot{X} = H_{21} + H_{22}X - XH_{11} - XH_{12}X \label{eq:riccati}
% \end{equation}
% First we see how the results in \ref{subsec:ricatti} apply to the Hamiltonian
Applying to the qudit Hamiltonian \ref{eq:rotatingframe} with $N_{\rm{max}}=2$,
the Riccati equation becomes
$$i\dot{x}_1 = (2\delta+\Delta) x_1 - \frac{1}{2}(\Omega_x + i\Omega_y)x_2 - \frac{\lambda}{2}(\Omega_x - i\Omega_y)x_1 x_2,$$
$$i\dot{x}_2 = \frac{\lambda}{2}(\Omega_x + i\Omega_y) + (2\delta+\Delta) x_2 - \frac{1}{2}(\Omega_x - i\Omega_y)x_1 - \frac{\lambda}{2}(\Omega_x - i\Omega_y)x_2^2.$$
If we define
\begin{equation}
  A = -\frac{1}{2}(x_2^2+\lambda x_1)\label{eq:A}
\end{equation}
\begin{equation}
  B = \frac{1}{2}x_1^2 \label{eq:B}
\end{equation}
\begin{equation}
  C = i(\dot{x_1}x_2-x_1\dot{x_2}) \label{eq:C}
\end{equation}
Then $\Omega_x$, $\Omega_y$, and $\delta(t)$ can be expressed
\begin{equation}
  \tilde{\Omega} = \Omega_x+i\Omega_y = \dfrac{A^\star C - B C^\star}{|A|^2-|B|^2}, \label{eq:riccati_Omega}
\end{equation}

\begin{equation}
  \delta(t) = \frac{1}{2} \left( \frac{i\dot{x}_1 + \frac{1}{2} x_2 \tilde{\Omega} + \frac{1}{2}\lambda x_1 x_2 \tilde{\Omega}^*}{x_1} - \Delta_2 \right) \label{eq:riccati_delta}.
\end{equation}



\subsubsection{Variational method}
One way to utilize \eref{eq:riccati_Omega} and \eref{eq:riccati_delta} is to determine $x_1$ and $x_2$ first and
solve numerically for $\tilde{\Omega}$ and $\delta$.
There's multiple difficulties with this approach. Since the Riccati equation \eref{riccati} is not linear, adjust the
pulse area of $\Omega_x$ is difficult. To overcome this issue, we first derive the and then
In this study, Riccati engineering method is utilized as an optimization method,
To preserve the X gate condition \eref{eq:xgate},
\begin{equation}
  \delta \Theta = \int_0^T \sum_{k=1}^2 \left( \frac{\partial \Omega_x}{\partial x_k} \delta x_k + \frac{\partial \Omega_x}{\partial \dot{x}_k} \delta \dot{x}_k + \text{c.c.} \right) dt
\end{equation}
must vanish. ($\dfrac{\partial}{\partial x_i}$ denotes Wirtinger derivatives.)
Integral by parts transforms the above equation into

\begin{equation}
  \delta \Theta = \left[ \sum_{k=1}^2 \left( P_k(T) \delta x_k(T) + P_k^*(T) \delta x_k^*(T) \right) \right] + \int_0^T \sum_{k=1}^2 \left( G_k(t) \delta x_k(t) + G_k^*(t) \delta x_k^*(t) \right) dt.
\end{equation}
Note that $\delta x_k(0) = 0$.


$$G_k(t) = \frac{\partial \Omega_x}{\partial x_k} - \frac{d}{dt} \left( \frac{\partial \Omega_x}{\partial \dot{x}_k} \right)$$

$$P_1(t) = \frac{\partial \Omega_x}{\partial \dot{x}_1} = i \frac{A^* - B^*}{2(\vert{}A\vert{}^2 - \vert{}B\vert{}^2)} x_2$$

$$P_2(t) = \frac{\partial \Omega_x}{\partial \dot{x}_2} = -i \frac{A^* - B^*}{2(\vert{}A\vert{}^2 - \vert{}B\vert{}^2)} x_1$$


\subsection{CZ gate}
In this section, for brevity, the length of pulses are assumed to be $N$, an odd number.
As seen in \ref{subsec:leakage}, the leakage error is approximately expressed as $\dfrac{1}{4}|G(\Delta)|^2$.
In the following we propose several methods to mitigate this error.
\subsubsection{anti-symmetric Chebyshev pulse}
Maintaining a specified sidelobe level, Chebyshev pulse is the signal that minimizes the mainlobe width

Mathematically, a chebyshev pulse $g_{ch}$ is a discrete time pulse defined as
$$
  g_{\text{ch}}[n] = \frac{1}{N} \Bigg[ 1 + 2r \sum_{k=0}^{(N-1)/2} (-1)^k T_{N-1}\left(x_0 \cos \frac{\pi k}{N}\right)  \cos \left( \frac{2\pi k}{L} \left(n + \frac{1}{2}\right) \right) \Bigg], $$
where $r = \frac{1}{T_{N-1}(x_0)}$ corresponds to the suppressed sidelobe level.
We call its anti-symmetrization
$$
  g_{ach}[n] = \begin{cases}
    g_{ch}[n]  \quad     & n< (N+1)/2  \\
    0             \quad  & n = (N+1)/2 \\
    -g_{ch}[N-n]   \quad & n > (N+1)/2 \\
  \end{cases}
$$
anti-symmetric Chebyshev pulse $g_{ach}$.
\subsubsection{Correctional Derivative pulse}
To reduce the fourier transform at $\Delta$, recall that
$$
  \mathcal{F}[f^\prime](\Delta) = i\omega \mathcal{F}[f](\Delta).
$$
Hence, by adding the second derivative to $f$, we can reduce the fourier transform to $0$,
$$
  \mathcal{F}[(1+D^2/\Delta^2)f](\Delta) = 0.
$$
This method is only applicable to pulses whose second derivative is anyti-symmetric, such as anti-symmetric Chebyshev,
to maintain the anti-symmetric property after adding the second derivative.


\subsubsection{Combining two pulses}
Combining two pulses of different origins
\subsubsection{Redefining the Slepian eigenvalue problem}
The (second) Slepian pulse is the solution of the second largest eigenvalue to the following problem
\begin{equation}
  \dfrac{\int_{-W}^{W}\,|G_{\text{DTFT}}(\omega)|^2 d\omega}
  {\int_{-\pi}^{\pi}\,|G_{\text{DTFT}}(\omega)|^2 d\omega} = \dfrac{g^TA_{m,n}g}{g^Tg} \label{eq:slepproblem}
\end{equation}
where $A_{m,n} = \sinc(\pi(m-n)W) = \dfrac{\sin(\pi(m-n)W)}{\pi(m-n)W}$.
(\ref{eq:slepproblem}) is the eigenvalue problem of the matrix $A_{m,n}$. It was Slepian who found a commuting
matrix for $A_{m,n}$ and solev teh problem analyitcally.

Close investigation of the \eref{eq:leakageerror}



For the targetted range of $W_1\leq \dfrac{T}{N}  \Delta \leq W_2$
$$\text{minimize } \frac{\int_{W_1}^{W_2} \vert{}G_{\text{DTFT}}(\omega)\vert{}^2 d\omega}{\int_{W_1}^{W_2} \vert{}G_{\text{DTFT}}(\omega)\vert{}^2 d\omega} = \frac{\tilde{g}^T (A_{\text{rect}}(W_2) - A_{\text{rect}}(W_1)) \tilde{g}}{\tilde{g}[0]^2 + \dots + \tilde{g}[N]^2}$$
To not elongate gate time we impose the following condition,
$\tilde{g}$ should be anti-symmetric and non-negative on the first half.
\begin{itemize}
  \item $\tilde{g}[s] = \tilde{g}[N-1-s]$
  \item $\tilde{g}[s] \geq 0 \quad \text{for }0\leq s\leq L-1$
\end{itemize}
Put
$$B = \begin{pmatrix} 1 & & \\ & \ddots & \\  &  & 1\\0 & \cdots & 0 \\  &  & -1 \\& \mathinner{\mkern2mu\raise1pt\hbox{.}\mkern2mu\raise4pt\hbox{.}\mkern2mu\raise7pt\hbox{.}\mkern1mu} & \\ -1 & &  \end{pmatrix}, \quad C = B^T A B$$
$$\begin{pmatrix} \tilde{g}[0] \\ \vdots \\ \tilde{g}[N] \end{pmatrix} = B \begin{pmatrix} h[0] \\ \vdots \\ h[L-1] \end{pmatrix}$$
Then
$$\text{maximize } -\frac{h^T C h}{h^T h}$$
is the problem to solve.
\subsubsection{Sequential Quadratic Programming}
Given an ansatz and an evaluation point $\Delta$, the algorithm goes as follows:

\begin{algorithm}
  \caption{Variational Quantum Circuit Optimization}\label{alg:vqa_opt}
  \begin{algorithmic}[1] % The [1] enables line numbering
    \Require Target unitary $U$, Error tolerance $\epsilon$
    \Ensure Optimized parameters $\boldsymbol{\theta}$
    \State Initialize $\boldsymbol{\theta} \gets \text{random values in } [0, 2\pi]$
    \State Set learning rate $\alpha \gets 0.01$
    \While{cost function $\mathcal{L}(\boldsymbol{\theta}) > \epsilon$}
    \State Compute gradient $\nabla \mathcal{L}(\boldsymbol{\theta})$ using parameter-shift rule
    \If{$\|\nabla \mathcal{L}(\boldsymbol{\theta})\| < \text{threshold}$}
    \State \textbf{break} \Comment{Early stopping criteria}
    \EndIf
    \State $\boldsymbol{\theta} \gets \boldsymbol{\theta} - \alpha \nabla \mathcal{L}(\boldsymbol{\theta})$
    \EndWhile
    \State \Return $\boldsymbol{\theta}$
  \end{algorithmic}
\end{algorithm}

